%%% FIELD shortest-normal-form-statement
There are \(G\in\mathfrak D(V)\), an absent directed edge \(e\) whose addition merges two directed SCCs, and a singleton E-separation query that holds in \(G\) and fails in \(G+e\), although every shortest \(D(B,C)\)-open witness for the failure in \(\operatorname{Lift}(G+e)\) contains no lifted occurrence generated by \(e\).  Hence path minimality alone does not yield a witness occurrence of \(e\) with \(e\)-free residual prefix and suffix; a normal form for this case must instead include the separate directed-ancestry certificate through which \(e\) activates a collider.
%%% FIELD shortest-normal-form-proof
Take \(V=\{a,b,c,d\}\).  Let the directed and bidirected edge sets of \(G\) be
\[
 E^{\to}(G)=\{(d,c)\},
 \qquad
 E^{\leftrightarrow}(G)=\bigl\{\{a,c\},\{c,b\}\bigr\},
\]
and put \(e=(c,d)\).  In the directed part of \(G\), the vertices \(c\) and \(d\) lie in different singleton SCCs, while \(d\to c\); after adding \(c\to d\), they are mutually reachable.  Thus \(e\) merges two SCCs.

Consider \(A=\{a\}\), \(B=\{b\}\), and \(C=\{d\}\).  The conditioning definition gives
\[
 D(B,C)=\{d_0,d_1\}.
\]
In \(\operatorname{Lift}(G)\), every edge incident to either copy of \(c\) has a head at that copy: the two original bidirected edges have heads at both ends, and every lifted arrow generated by \(d\to c\) has its head at a copy of \(c\).  Every walk from \(a_0\) to \(b_1\) must pass through a copy of \(c\) as an interior occurrence, because \(a\) is adjacent only to \(c\) and \(b\) is adjacent only to \(c\).  At the first such passage between the \(a\)-side and the \(b\)-side, the occurrence of \(c\) is a collider.  No copy of \(c\) is a directed ancestor of \(D(B,C)\) in \(\operatorname{Lift}(G)\): the only directed edge points from \(d\) toward \(c\), and there is no directed edge out of \(c\).  The collider-activation rule therefore blocks every such walk.  Hence
\[
 (\{a\},\{b\},\{d\})\in\mathcal I_E(G).
\]

In \(\operatorname{Lift}(G+e)\), the walk
\[
 a_0\leftrightarrow c_0\leftrightarrow b_1
\]
is open given \(D(B,C)\).  Its only interior occurrence is the collider \(c_0\), and the new lifted arrow \(c_0\to d_0\) makes \(c_0\) an ancestor of the conditioned vertex \(d_0\).  Both endpoints are outside \(D(B,C)\).  Thus the query fails in \(G+e\).

There is no one-edge walk from \(a_0\) to \(b_1\), because \(G+e\) has no original edge between \(a\) and \(b\).  The displayed two-edge walk is consequently shortest.  Moreover, every two-edge walk from \(a_0\) to \(b_1\) has a middle copy of \(c\): no lifted edge joins a copy of \(a\) to a copy of \(d\), or a copy of \(d\) to \(b_1\).  Its two edges are therefore lifted bidirected edges generated by \(\{a,c\}\) and \(\{c,b\}\), not lifted occurrences of \(e\).  Hence every shortest newly open witness avoids \(e\).  The addition acts only through the directed-ancestry certificate that activates the old collider.  This proves that the requested occurrence cannot be extracted from witness minimality alone; it must be sought in an augmented witness-certificate object.
%%% FIELD marked-probe-statement
There are a three-vertex set \(V=\{s,u,v\}\), graphs \(G,H\in\mathfrak D(V)\), and a directed edge \(e=(u,v)\in E(H)\setminus E(G)\) that merges two SCCs when added to \(G\), such that \(G\equiv_E H\), but the following marked prefix boundary differs.  For the fixed lifted occurrence \(u_0\to v_0\) of \(e\), \(\operatorname{Lift}(G)\) contains an \(e\)-free prefix from \(s_0\) to \(u_0\) whose final incidence at \(u_0\) is a tail, whereas \(\operatorname{Lift}(H)\), with all lifted occurrences of \(e\) forbidden, has no walk ending at \(u_0\) with a tail at its final incidence.  Consequently this marked prefix signature is not determined by any collection of ordinary E-separation queries, even with \(B=\{s\}\), \(C=\varnothing\), and hence \(D(B,C)=\varnothing\), held fixed.
%%% FIELD marked-probe-proof
Let \(V=\{s,u,v\}\), and give both graphs complete bidirected support, including all three bidirected loops:
\[
 E^{\leftrightarrow}(G)=E^{\leftrightarrow}(H)=U_{\leftrightarrow}.
\]
Define their directed parts by
\[
 E^{\to}(G)=\{(v,u),(u,s)\},
 \qquad E^{\to}(H)=\{(v,u),(u,v)\},
 \qquad e=(u,v).
\]
Thus \(e\in E(H)\setminus E(G)\).  We first prove \(G\equiv_EH\), rather than assuming it.  Fix \(a,b\in V\) and \(C\subseteq V\).  By the lifted-conditioning definition,
\[
 D(\{b\},C)=C_0\cup\bigl(C_1\setminus\{b_1\}\bigr).
\]
If \(a\in C\), then \(a_0\in D(\{b\},C)\), so the endpoint condition in the definition of a \(\sigma\)-open walk makes the singleton query separated in both graphs.  If \(a\notin C\), then neither endpoint \(a_0,b_1\) is in \(D(\{b\},C)\).  Complete bidirected support and the lift definition give the lifted edge \(a_0\leftrightarrow b_1\): when \(a=b\) it is the cross-copy edge generated by \(\{a,a\}\), and when \(a\ne b\) it is generated by \(\{a,b\}\).  This one-edge walk has no interior occurrences and is therefore \(\sigma\)-open.  Thus \(\Sigma_E(G)=\Sigma_E(H)\).  The singleton-composition lemma now gives \(\mathcal I_E(G)=\mathcal I_E(H)\), which is \(G\equiv_EH\) by the equivalence definition.

The directed part of \(G\) has distinct singleton SCCs \(\{u\}\) and \(\{v\}\): it contains \(v\to u\), but its only arrow out of \(u\) goes to \(s\), from which no directed arrow leaves.  Adding \(e=u\to v\) makes \(u\) and \(v\) mutually reachable, so \(e\) changes the SCC partition.  By the lift definition, \(e\) generates the fixed occurrence \(u_0\to v_0\).  Also, \(u\to s\) in \(G\) generates \(u_0\to s_0\); traversed from \(s_0\) to \(u_0\), it is a one-edge prefix whose final mark at \(u_0\) is a tail.  In \(G+e\), following this prefix by \(u_0\to v_0\) produces the concrete tail--tail incidence pattern at the left boundary of the selected lifted occurrence.

In \(H\), forbid all three lifted arrows generated by \(e\).  Every remaining lifted edge incident to \(u_0\) is either bidirected, and hence has a head at \(u_0\), or is \(v_0\to u_0\), which also has a head at \(u_0\).  There is no directed loop at \(u\), and \(H\) omits \(u\to s\).  Hence no remaining walk can end at \(u_0\) with a tail at its final incidence.  This conclusion holds in particular for the fixed choice \(B=\{s\}\), \(C=\varnothing\), for which the conditioning definition gives \(D(B,C)=\varnothing\).

The two graphs agree on their entire E-model but disagree on this marked property.  Therefore the property cannot be a function of all ordinary E-separation answers; a fortiori it cannot be decided by finitely many such answers with any fixed \(B,C\).  This disproves the marked-probe completeness sublemma required by the proposed unrestricted splice, while making no claim that unrestricted single-edge transfer itself is false.
%%% FIELD single-edge-proof
Fix \(G\in\mathfrak D(V)\) and \(e\in E(\Gamma(G))\setminus E(G)\).  The fixed-profile proposition gives
\[
 \mathcal I_E(\Gamma(G))=\mathcal I_E(G),
 \qquad\text{hence}\qquad \Gamma(G)\equiv_EG.
\]
The edge assumption gives the componentwise inclusions
\[
 G\subseteq G+e\subseteq\Gamma(G).
\]
Applying the interval conclusion of the antitonicity lemma to these two displays yields \(G+e\equiv_EG\).

The profile-saturation definition identifies the possible new \(e\)'s exactly: a directed edge internal to a non-singleton directed SCC, a directed edge from a vertex whose incoming support to that SCC is already occupied, or a bidirected edge in an already occupied rectangle between two distinct SCCs.  The same definition inserts directed loops only in non-singleton SCCs and inserts neither internal bidirected edges nor bidirected loops.  Thus the proof covers precisely the stated cases.  Finiteness of \(V\) is used only to keep all graphs in \(\mathfrak D(V)\); no positivity, rank, support, overlap, boundary, or limit condition is involved.
%%% FIELD greatest-proof
The fixed-profile proposition directly gives, for each \(G\in\mathfrak D(V)\),
\[
 \Gamma(G)\in\mathfrak F_G,
 \qquad \Gamma(G)\equiv_EG,
\]
closure of \(\mathfrak F_G\) under componentwise union, and the fact that every \(H\in\mathfrak F_G\) satisfies \(H\subseteq\Gamma(G)\).  If \(L\in\mathfrak F_G\) is also greatest, then \(L\subseteq\Gamma(G)\) and \(\Gamma(G)\subseteq L\), so antisymmetry of edge inclusion gives \(L=\Gamma(G)\).  This proves the unconditional fixed-profile assertion.

For the unrestricted characterization, let \(\mathsf T\) denote unrestricted single-edge transfer, let \(\mathsf U\) denote binary-union closure of every E-equivalence class, and let \(\mathsf M\) denote existence of a greatest member in every E-equivalence class.

Assume \(\mathsf T\), and take \(X\equiv_EY\).  Because the edge universe \(U_{\to}\cup U_{\leftrightarrow}\) is finite, write
\[
 E(Y)\setminus E(X)=\{e_1,\ldots,e_m\},\qquad
 K_0=X,\qquad K_r=K_{r-1}+e_r.
\]
Inductively, \(K_{r-1}\equiv_EX\equiv_EY\) and \(e_r\in E(Y)\setminus E(K_{r-1})\).  Property \(\mathsf T\) therefore gives \(K_r\equiv_EK_{r-1}\).  At \(r=m\), \(K_m=X\cup Y\equiv_EX\), proving \(\mathsf T\Rightarrow\mathsf U\).

Assume \(\mathsf U\).  The class \(\mathfrak D(V)\) is finite because both edge universes are finite.  Hence each nonempty E-equivalence class can be enumerated as \(X_1,\ldots,X_N\).  Repeated application of \(\mathsf U\) shows that \(X_1\cup\cdots\cup X_N\) remains in the class.  It contains every class member and is therefore greatest; uniqueness follows from antisymmetry.  Thus \(\mathsf U\Rightarrow\mathsf M\).

Assume \(\mathsf M\), take \(X\equiv_EY\), and let \(L\) be the greatest member of their class.  Then
\[
 X\subseteq X\cup Y\subseteq L,
 \qquad X\equiv_EL.
\]
The interval conclusion of the antitonicity lemma gives \(X\cup Y\equiv_EX\), so \(\mathsf M\Rightarrow\mathsf U\).  Finally, under \(\mathsf U\), if \(X\equiv_EY\) and \(e\in E(Y)\setminus E(X)\), then
\[
 X\subseteq X+e\subseteq X\cup Y,
 \qquad X\equiv_E X\cup Y.
\]
The same interval lemma gives \(X+e\equiv_EX\), proving \(\mathsf U\Rightarrow\mathsf T\).  Hence \(\mathsf T\), \(\mathsf U\), and \(\mathsf M\) are equivalent.  This is a characterization, not an assertion that they hold outside the fixed-profile families.
%%% FIELD saturation-proof
Define the local abbreviation
\[
 \operatorname{Adm}_{\Gamma}(G)
 :=\operatorname{Adm}(G)\cap E(\Gamma(G)).
\]
By the admissibility definition, every member of \(\operatorname{Adm}(G)\) is absent from \(G\); hence
\[
 \operatorname{Adm}_{\Gamma}(G)
 \subseteq E(\Gamma(G))\setminus E(G).
\]
Conversely, let \(e\in E(\Gamma(G))\setminus E(G)\).  The proved fixed-profile single-edge result gives \(G+e\equiv_EG\), so their E-models, and therefore their singleton signatures, are equal.  The admissibility definition gives \(e\in\operatorname{Adm}(G)\).  Thus
\[
 \operatorname{Adm}_{\Gamma}(G)=E(\Gamma(G))\setminus E(G).
\]
Adding both sides to \(G\), and then invoking the fixed-profile proposition, proves
\[
 G+\operatorname{Adm}_{\Gamma}(G)=\Gamma(G),
 \qquad
 \mathcal I_E\bigl(G+\operatorname{Adm}_{\Gamma}(G)\bigr)=\mathcal I_E(G).
\]

Now suppose, only for the conditional unrestricted conclusion, that \(G^\star\equiv_EG\).  If \(e\in\operatorname{Adm}(G)\), then \(\Sigma_E(G+e)=\Sigma_E(G)\) by definition; singleton composition gives \(G+e\equiv_EG\).  Since \(G^\star\) is the componentwise union of all equivalent graphs, \(e\in E(G^\star)\setminus E(G)\).  Conversely, if \(e\in E(G^\star)\setminus E(G)\), then
\[
 G\subseteq G+e\subseteq G^\star,
 \qquad G\equiv_EG^\star.
\]
The interval conclusion of antitonicity gives \(G+e\equiv_EG\), hence equality of singleton signatures and \(e\in\operatorname{Adm}(G)\).  We have proved
\[
 \operatorname{Adm}(G)=E(G^\star)\setminus E(G),
 \qquad
 G+\operatorname{Adm}(G)=G^\star,
 \qquad
 \mathcal I_E\bigl(G+\operatorname{Adm}(G)\bigr)=\mathcal I_E(G)
\]
under exactly the displayed condition.  All arguments are finite and graphical; no positivity, rank, support, overlap, boundary, or limit premise is used.
%%% FIELD oracle-proof
For each ordered \((a,b)\in V^2\), there are \(2^{n-1}\) subsets \(C\subseteq V\setminus\{a\}\).  The number of queried entries is therefore
\[
 n^2 2^{n-1}=Q.
\]
Consider an omitted singleton entry, so \(a\in C\).  The conditioning definition gives \(a_0\in C_0\subseteq D(\{b\},C)\); the endpoint condition for \(\sigma\)-openness therefore makes this query separated.  Consequently the extension rule in the corrected recovery definition reconstructs the entire singleton signature \(\Sigma_E(G)\).  By the singleton-composition lemma, this signature determines the full E-model \(\mathcal I_E(G)\).

The sets \(U_{\to}\) and \(U_{\leftrightarrow}\) are finite because \(V\) is finite, so \(\mathfrak D(V)\) is finite.  Exhaustive comparison with the reconstructed signature therefore terminates and returns exactly
\[
 \mathcal C(s)=\{H\in\mathfrak D(V):H\equiv_EG\}.
\]
This set is nonempty because it contains \(G\).  The first component returned by \(\operatorname{Alg}_E\) is exactly the componentwise union in the definition of \(G^\star\), and its second component is exactly the edge intersection in the definition of \(\operatorname{Comp}(G)\).  Hence, without a closure assumption,
\[
 \operatorname{Alg}_E(s)=\bigl(G^\star,\operatorname{Comp}(G)\bigr).
\]
This equality identifies an observable union functional and an observable intersection functional; it does not by itself place the union back in the E-equivalence class.

Assume now the additional displayed condition \(G^\star\equiv_EG\).  Every equivalent graph is contained in \(G^\star\) by its union definition, so \(G^\star\) is an in-class greatest member and is unique by antisymmetry.  Fix \(e\in E(G^\star)\).  If
\[
 \Sigma_E(G^\star-e)=\Sigma_E(G^\star),
\]
singleton composition gives \(G^\star-e\equiv_EG^\star\).  This equivalent graph omits \(e\), so the intersection definition of compelled edges gives \(e\notin\operatorname{Comp}(G)\).

Conversely, if \(e\notin\operatorname{Comp}(G)\), the intersection definition supplies \(K\equiv_EG\) with \(e\notin E(K)\).  Greatestness and the omission of \(e\) give
\[
 K\subseteq G^\star-e\subseteq G^\star,
 \qquad K\equiv_EG^\star.
\]
The interval conclusion of the antitonicity lemma yields \(G^\star-e\equiv_EG^\star\), hence equal singleton signatures.  Taking contrapositives in both directions proves
\[
 e\in\operatorname{Comp}(G)
 \quad\Longleftrightarrow\quad
 \Sigma_E(G^\star-e)\ne\Sigma_E(G^\star)
\]
under \(G^\star\equiv_EG\).  There are no statistical positivity, rank, support, overlap, boundary, or limiting assumptions: this is exact finite graphical identification from a promised oracle.
